% Input chapter: deterministic definition and lift removal.
\section{The structural profile}
\label{profile:definition}

Throughout, $k\geq6$ is fixed, $H=k+1$, and $V_i=\sum_{g\leq i}L_g$,
with $V_0=0$. We first define a numerical function on $L_i\geq1$.
The final definition will cover $L_i\geq0$. The optimization dimension
and every finite enumeration below depend only on $k$.

The four factors in the profile describe four constraints on the same
family of decompositions: exponential tilting, proper exit counts,
conditional binary survival, and the terminal ordinary bridge. We define
them in this order. Although the final expression is a product, the
underlying events share counts and environments; their joint realization
and their simultaneous necessity are proved separately below.

\subsection{The saddle and the original named groups}
On the closed ordered simplex $1\geq s_1\geq\cdots\geq s_k\geq0$ put
\begin{align}
 G_k&=s_k\log2,&
 G_i&=s_i\log(1+\exp(G_{i+1}/s_i)),\\
 Z_i&=\exp G_i,&
 J(L;s)&=\sum_i L_i+\sum_j s_jV_j-\sum_i L_iZ_i.
 \label{profile:saddle}
\end{align}
At zero products use the continuous perspective extension. Write $s=s(L)$ for the
unique maximizer, whose existence and uniqueness are proved in
Section~\ref{paper:geometry}. The geometric theorem supplies
\begin{equation}
 \tfrac12<s_k\leq s_1<1-\epsilon_H,
 \qquad \epsilon_H=\frac{2\log2-1}{2H(\log H)^2}.
 \label{profile:compact}
\end{equation}
In particular $1\leq Z_i\leq H$. A maximal consecutive group
$[a,b]$ of equal products is called an owner; write its product as
$\sigma$. A proper owner has $b<k$ and
$z=\exp(G_{b+1}/\sigma)$; the terminal owner has $b=k$ and $z=1$.
The exterior names of this owner are $a,\ldots,b$. Its native groups
are the original groups $i\in[a,b]$, of ranks $r_i=b-i+1$.
Original groups $g<a$ are its old groups, of rank $b-a+1$.
Each original group retains its own random clock.

In the divisor model, name $c$ is coordinate $d_c$: primes in group $i$
may be assigned to $c\geq i$, and $V_c$ is its deadline. Assignments to
coordinates beyond $b$, together with the unused primes, form the
owner's compound child.

For $l<k$ define
\[
 \beta_l=\frac{\exp(G_{l+1}/s_l)}{1+\exp(G_{l+1}/s_l)}.
\]
The reference intensity per eligible exterior name is
\begin{equation}
 \kappa_g=
 \begin{cases}
 Z_g\displaystyle\prod_{l=g}^{a-1}\beta_l/(b-a+1+z),&g<a,\\
 (b-g+1+z)^{\sigma-1},&a\leq g\leq b.
 \end{cases}
 \label{profile:old-coefficients}
\end{equation}
Each old group retains the intensity acquired along its original route
into this owner. Averaging these coefficients would change the
constraints at the individual deadlines. Set
\[
 E_q(z)=(q+z)\log(q+z)-z\log z,\qquad
 g_g(q)=q-\kappa_gE_q(z).
\]

\subsection{An explicit enumeration of the two endpoint scores}
For every native group $i$ enumerate the edges $r_i\to q$, where
$0\leq q<r_i$ for a proper owner, and $1\leq q<r_i$ for the terminal
owner. The latter exclusion removes the degenerate ordinary edge.
The retained exterior names are $B=\{b-q+1,\ldots,b\}$; $B$ is empty
when $q=0$. The compound child is always retained.

Here is the complete named outside-context enumeration for that edge.
For each name $c\in\{a,\ldots,i-1\}$ choose
$\tau_c\in\{V_0,\ldots,V_c\}$ independently. For each $c\in B$ choose
$\tau_c\in\{V_i,\ldots,V_c\}$ independently. In the left completion
set $\tau_c=V_{i-1}$ for all $i\leq c\leq b-q$; in the right completion
set those same cuts to $V_i$. All other cuts are identical in the two
completions. This gives a finite set of lawful pairs $(v^-,v^+)$.
Every cut respects its original name's deadline, and the compound child
remains intact. These choices specify the entire enumeration.

For either completion $v$ and original group $1\leq g\leq b$, let
\[
 A_{v,g}=\{c\in\{a,\ldots,b\}:\tau_c\geq V_g\},
 \qquad q_{v,g}=|A_{v,g}|,
 \qquad t_v=\max_{a\leq c\leq b}\tau_c.
\]
In an old group the shallowest eligible name is $a$, and in native
group $g$ it is $g$. Call $g$ good when $0<q_{v,g}<r_g$ and this
shallowest eligible name is absent from $A_{v,g}$. Define
\begin{equation}
 \overline C_v=\sum_{g\leq b}L_gg_g(q_{v,g}),\qquad
 S_{\mathrm{good}}(v)=\sum_{g\text{ good}}\sqrt{L_g}.
 \label{profile:reference}
\end{equation}
The endpoint geometry gives $\overline C_v\geq0$ at the physical
optimizer. Both the old coefficients and the old lengths are included.

For a proper owner put $T=\sum_{a\leq i\leq b}L_i$,
$d=\sigma/s_{b+1}-1$, and
\[
 \widehat R=1+\frac{\sqrt{1+T}}{1+d\sqrt{1+T}}.
\]
For the terminal owner set $\widehat R=0$. Define
\begin{equation}
 S_v=\begin{cases}
 1+\widehat R+\overline C_v,&t_v\leq V_b/2,\\
 1+\overline C_v+S_{\mathrm{good}}(v),&t_v>V_b/2,
 \end{cases}
 \qquad X_e=\min_v S_{v^-},\quad Y_e=\min_v S_{v^+}.
 \label{profile:scores}
\end{equation}
The early score includes the uncertainty of a proper exit-count window.
At the late endpoint the incoming contribution cancels from the intrinsic
cost; the additional square-root terms instead measure the fluctuation
needed to remove an eligible shallow name while retaining deeper names.
The two minima are taken separately over the same finite set of paired
lawful contexts. The lower bound makes the two scores compatible by
keeping the actual type increment $\Delta$ fixed and choosing the legal pair
$x=\max(a_*X_e,a_*Y_e-\Delta)$, $y=x+\Delta$; see
\eqref{lower:local-pair}.

\subsection{The independent scalar and all canonical edge factors}
Let $E,W$ be independent stationary standard Ornstein--Uhlenbeck
processes, with covariance $e^{-|u-v|/2}$. For $0<\eta,\rho<1$ set
\begin{align}
 F_{\eta,\rho}(t)&=
 \mathbb E_E\Big[\mathbb P_W\{
 \sqrt\rho E_u+\sqrt{1-\rho}W_u\geq0\ (0\leq u\leq t)
 \mid E\}^{\eta}\Big],\\
 \chi(\eta,\rho)&=\lim_{t\to\infty}-t^{-1}\log F_{\eta,\rho}(t).
 \label{profile:chi}
\end{align}
The scalar theorem proves existence, local Lipschitz continuity, and
$\eta/2\leq\chi\leq1/2$. It is independent of $K_k$ and of all
multirank history spaces. It also gives certified computation: if
$b(t)=-\log F_{\eta,\rho}(t)$ and
$p_g=(1+e^{-g/2})/(1-e^{-g/2})$, then
\begin{equation}
 \frac{b(t)}{p_g(t+g)}\leq\chi(\eta,\rho)\leq\frac{b(t)}t.
 \label{profile:chi-bracket}
\end{equation}
The scalar finite Gaussian quadrature and crossing error bounds make
each endpoint effective. Taking $g=2\log t$ gives vanishing width.

For $e=(i,r,q)$ put $a=\log(r+z)$, $b_e=\log(q+z)$,
$p=(q+z)/(r+z)$, and
\begin{equation}
 m=a-b_ep,\qquad
 \rho_e=\frac{m^2}{m^2+b_e^2p(1-p)},\qquad
 T_e=\frac{(1+X_e)(1+Y_e)}{1+(r+z)^\sigma L_i},\qquad
 P_e=\min(1,T_e)^{2\chi(\sigma,\rho_e)}.
 \label{profile:edge-factor}
\end{equation}
These laws are nondegenerate. In particular a proper zero-output edge
has $z>1$, whereas the terminal zero-output edge was excluded. Every
canonical edge is used, including edges with a bounded positive factor.

\subsection{An explicit ordinary suffix}
Here $h_i=k-i+2$ on the complete original chain. If the selected suffix
starts at group $j$, then $V_{\rm pre}=V_{j-1}$ includes all preceding
owners as well as the earlier terminal groups.
For $2\leq h\leq H$ define
\[
 \alpha_h=\frac{\log((h-1)/\log h)}{\log h},\quad
 c_w=\frac1{4H^2},\quad
 \tau=\min\left\{\frac14,\frac{c_w}{2H(\log H)^2},
 \frac14\min_{2\leq h<H}(\alpha_{h+1}-\alpha_h)\right\}.
\]
The thresholds $\alpha_h$ are strictly increasing. If the terminal
product $\sigma$ lies in $[\alpha_{h_i}-\tau,\alpha_{h_i}+\tau]$ for
a terminal native group, begin the suffix at that group; at most one
qualifies. Otherwise begin at the first terminal native group with
$e_i=h_i^\sigma\log h_i-(h_i-1)>0$. The KKT theorem ensures existence.
Let $U$ be the suffix length, $M$ its largest original block, and $A,B$
the total lengths before and after that block within the suffix. Break
largest-block ties at the earliest index. Put
\begin{equation}
 D=\sum_{i\text{ in suffix}}e_iL_i,\qquad
 \mathcal B(L)=\frac1{\sqrt{1+U}}
       \min\left\{1,\frac{(1+B)(1+D+A)}{1+M}\right\}.
 \label{profile:ordinary}
\end{equation}
Here $D\geq0$, and the suffix theorem gives $D\geq c_kV_{\rm pre}$
and $U\geq c_kV_k$. This is the sole ordinary root factor.

The positive-length raw expression is
\begin{equation}
 \Psi_k(L)=e^{-J(L;s(L))}\mathcal B(L)
 \prod_{j\text{ proper}}(1+d_j\sqrt{1+T_j})^{-\sigma_j}
 \prod_{e\text{ canonical}}P_e.
 \label{profile:raw}
\end{equation}
All lengths, reference costs and products in this expression refer to
one and the same physical parameter vector. The rate $J$ is a Poisson
entropy, and each persistence exponent splits into an environmental
selection cost and a conditional path cost. These interpretations,
proved in Section~\ref{fp:sec:interpretation}, also locate, for sufficiently
large positive lengths, a family of prime tuples contributing the full
order of the arithmetic sum.
\begin{figure}[htbp]
\centering
\begin{tikzpicture}[
  every node/.style={font=\small,align=center},
  box/.style={draw,rounded corners=2pt,inner sep=7pt},
  arr/.style={-{Stealth[length=2mm]},thin}]
\node[box,text width=10.8cm] (saddle)
 {One physical vector $L$ and its ordered saddle $s(L)$\\
  rate $J$, exact equality blocks, and original transported intensities};
\node[box,text width=3.05cm,below=9mm of saddle,xshift=-3.85cm] (count)
 {Proper exit windows\\[2pt]
  $(1+d\sqrt{1+T})^{-\sigma}$};
\node[box,text width=3.05cm,below=9mm of saddle] (binary)
 {Paired endpoint scores\\[2pt]
  $\min(1,T_e)^{2\chi(\sigma,\rho_e)}$};
\node[box,text width=3.05cm,below=9mm of saddle,xshift=3.85cm] (ordinary)
 {One terminal suffix\\[2pt]
  $\mathcal B(L)$};
\draw[arr] (saddle.south -| count.north) -- (count.north);
\draw[arr] (saddle.south) -- (binary.north);
\draw[arr] (saddle.south -| ordinary.north) -- (ordinary.north);
\node[box,text width=10.8cm,below=9mm of binary] (profile)
 {$\Psi_k(L)=e^{-J(L)}\,
    \displaystyle\prod_{O\text{ proper}}\text{exit factor},
    \displaystyle\prod_{e\text{ canonical}}\text{binary factor},
    \mathcal B(L)$};
\draw[arr] (count.south) -- (profile.north -| count.south);
\draw[arr] (binary.south) -- (profile.north);
\draw[arr] (ordinary.south) -- (profile.north -| ordinary.south);
\end{tikzpicture}
\caption{The deterministic construction of the positive-length profile.
The arrows indicate dependence on the same saddle data. The product
formula is established by the lower and upper proofs; it does not
assert independence of the count and path events.}
\label{fig:profile-structure}
\end{figure}


\subsection{The profile on the full parameter domain}
For $\ell_i\geq\log3$ put $L_i=\ell_i-\log3$, $A_k=2(k+1)^3$,
and define
\begin{equation}
 \Phi_k(\ell)=\inf_{\lambda\geq1}
     e^{A_k\lambda}\Psi_k\big((\max(L_i,\lambda))_{i=1}^k\big).
 \label{profile:final}
\end{equation}
This completes the numerical definition with one scalar optimization;
$\chi$ is the universal scalar specified in \eqref{profile:chi}.


Lemma~\ref{profile:finite-search} gives an explicit finite interval for
this last infimum and proves that $0<\Phi_k(\ell)<\infty$.
Every candidate length vector uses its own maximizer, exact equality
blocks, endpoint contexts and ordinary suffix.
